% 第27章思考题27.1--27.7的独立概念解析，候选稿，未获专家认证。
\section*{第27章\quad 薛定谔方程：思考题解析}
\paragraph{思考题 27.1}
薛定谔方程并非从经典牛顿力学\emph{严格逻辑必然地}
唯一推出，而是受能量—动量关系及德布罗意物质波
启发提出、由实验支持的量子运动基本假设。
非相对论单粒子经典Hamilton函数
\(H=p^2/(2m)+V(\boldsymbol r,t)\)，
对应算符替换
\(E\mapsto i\hbar\partial_t\)、
\(\boldsymbol p\mapsto-i\hbar\nabla\)，
从而得到
\[
\boxed{i\hbar\partial_t\Psi=
\left[-\frac{\hbar^2}{2m}\nabla^2+V\right]\Psi}.
\]
这只是启发式构造，数学上的自伴性、
定义域和实验验证不可省略。

\paragraph{思考题 27.2}
若Hamilton算符 \(H=H^\dagger\)，
按复共轭方程得
\(\partial_t|\Psi|^2=
-(\nabla\cdot\boldsymbol j)\)，其中
\[
\boxed{\rho=|\Psi|^2,\quad
\boldsymbol j=\frac{\hbar}{2mi}
(\Psi^*\nabla\Psi-\Psi\nabla\Psi^*)},
\]
\[
\boxed{\partial_t\rho+\nabla\cdot\boldsymbol j=0}.
\]
对整个空间积分且无边界净流：
\(d\int|\Psi|^2d^3r/dt=0\)。
因此初始归一化为1的波函数保持归一化。
若有非Hermitian有效Hamilton量或吸收性边界，
单粒子概率范数可能变化，需另行解释。

\paragraph{思考题 27.3}
标准一维非奇异有限势情形，
物理上可接受的定态波函数须单值、有限、
平方可积、可归一化，并满足边界条件；
通常 \(\psi,\psi'\) 连续。
但是无限深势壁使有限侧波函数在壁处为零，
\emph{一阶导数不必在无穷势阶跃处连续}；
\(\delta\) 势会让导数存在规定跳跃，
散射态的广义本征函数也可以不属于普通
\(L^2(\mathbb R)\)。
因此
\[
\boxed{\int|\psi|^2d^3r=1}
\]
是束缚概率态的归一化要求，不能机械施用于
所有广义散射平面波。

\paragraph{思考题 27.4}
Born概率解释为
\(P(\Omega)=\int_\Omega|\Psi|^2dV\)
（若\(\|\Psi\|_2=1\)）。
任意非零常数 \(C\) 乘以未归一化解
不会改变归一化后的概率比：
\[
P(\Omega)=
\frac{\int_\Omega|C\Psi|^2dV}
{\int_{\mathbb R^3}|C\Psi|^2dV}
=\frac{\int_\Omega|\Psi|^2dV}
{\int_{\mathbb R^3}|\Psi|^2dV}.
\]
唯有确定范数后，才能把 \(|\Psi|^2dV\)
直接当作概率，而不是相对概率密度。

\paragraph{思考题 27.5}
对\emph{一维、自伴、束缚、非简并}的
Sturm--Liouville体系，按能量递增编号，
基态内部无节点，第 \(n\) 个本征态有
\(\boxed{n-1}\) 个\emph{内部节点}
（约定基态编号 \(n=1\)）。
因此基态密度在允许区域的内部一般不为零，
激发态有零密度的节点 \(n-1\) 个；
无限壁处零点不计作内部节点。
这一规律并不能无条件推广到
二维或三维\emph{简并}态，也不能直接用于
任意叠加态。原图27.3/27.5/27.12
需逐图核对势壁、节点位置及禁阻区尾部。

\paragraph{思考题 27.6}
按图27.13阶梯左区 \(V_L<V_R\)，
同一能级 \(E_5>V_R\)。
在局部近似下
\[
k_j=\frac{\sqrt{2m(E_5-V_j)}}{\hbar},
\quad \lambda_j=\frac{2\pi}{k_j},\quad
|\psi_j|_{\rm WKB}\propto k_j^{-1/2}.
\]
因为右区 \(V_R\) 较高，
\(k_R<k_L\)，故
\[
\boxed{\lambda_R>\lambda_L},\qquad
\boxed{A_R>A_L\quad\text{(WKB平均振幅)}}.
\]
严格常阶梯势阱中局部波为不同波数的正弦组合，
需同时满足壁零值与阶梯处 \(\psi,\psi'\) 连续；
振幅的确切大小由全局边界条件确定，
上述振幅不等式不是任意薄小区间的严格定理。

\paragraph{思考题 27.7}
孤立、时间独立且自伴的Hamilton量
对每个能量本征态
\(\Psi_n(t)=\psi_ne^{-iE_nt/\hbar}\)，
其概率密度不随时间变，
\[
\boxed{P_n(t)=|\Psi_n(t)|^2=|\psi_n|^2}.
\]
它\emph{不会}在孤立模型中自发跃迁至其他本征态，
因此此模型的严格定态寿命可视为无限大。
真实原子激发态与量子化电磁场耦合，
完整体系中的原子子系统可通过自发辐射衰变；
有限寿命和能级自然宽度因此来自\emph{开放子系统}，
并不反驳孤立Hamilton定态方程。
